\documentclass[11pt,a4paper]{article}

\usepackage[utf8]{inputenc}
\usepackage[T1]{fontenc}
\usepackage[brazilian]{babel}
\usepackage{geometry}
\usepackage{graphicx}
\usepackage{booktabs}
\usepackage{float}
\usepackage{enumitem}
\usepackage{amsmath}
\usepackage{xcolor}
\usepackage{fancyhdr}
\usepackage{setspace}
\usepackage{hyperref}
\usepackage{tcolorbox}

\geometry{left=2.5cm,right=2.5cm,top=2.5cm,bottom=2.5cm}
\setlength{\parindent}{0pt}
\setlength{\parskip}{0.55em}
\onehalfspacing

\tcbuselibrary{skins,breakable}
\newtcolorbox{instrucao}[1][]{colback=blue!5,colframe=blue!40!black,
  fonttitle=\bfseries, title=#1, breakable}
\newtcolorbox{atencao}[1][]{colback=orange!10,colframe=orange!60!black,
  fonttitle=\bfseries, title=#1, breakable}

\pagestyle{fancy}
\fancyhf{}
\fancyhead[L]{\small FGA0092 -- Princípios de Comunicação para Engenharia}
\fancyhead[R]{\small Enunciado -- Prática 04}
\fancyfoot[C]{\small\thepage}
\renewcommand{\headrulewidth}{0.4pt}

\begin{document}

\begin{center}
  \includegraphics[width=0.18\textwidth]{../../figures_global/unb.png}\\[0.5cm]
  {\large\textbf{Universidade de Brasília (UnB) -- FCTE}}\\[0.2cm]
  {\normalsize FGA0092 -- Princípios de Comunicação para Engenharia}\\[0.4cm]
  \rule{14cm}{0.8pt}\\[0.3cm]
  {\LARGE\textbf{Prática 04}}\\[0.15cm]
  {\Large\textbf{Modulação SSB (Banda Lateral Única)}}\\[0.3cm]
  \rule{14cm}{0.8pt}
\end{center}

\vspace{0.3cm}
\textbf{Duração:} 50 minutos \hfill
\textbf{Software:} GNU Radio Companion $\geq$ 3.10 \hfill
\textbf{Semestre:} 2026/2

\bigskip

% ======================================================
\section{Objetivos}

\begin{itemize}
  \item Gerar o sinal DSB-SC e identificar as bandas laterais superior
    e inferior no espectro.
  \item Eliminar uma banda lateral usando filtro passa-banda, obtendo o
    sinal SSB.
  \item Comparar a largura de banda do DSB-SC com a do SSB no espectro.
  \item Demodular o sinal SSB por detecção coerente, medir o fator de escala
    e compará-lo com a mensagem original.
  \item Medir experimentalmente o efeito de um erro de frequência $\Delta f$
    no oscilador local.
  \item Gerar SSB pela transformada de Hilbert e comparar a supressão da banda
    lateral indesejada obtida pelos dois métodos.
\end{itemize}

% ======================================================
\section{Fundamentos teóricos}

O sinal DSB-SC ocupa largura de banda $B_\text{DSB} = 2 f_m$. O sinal SSB
elimina uma das bandas laterais e reduz a largura de banda pela metade:
\begin{equation}
  B_\text{SSB} = f_m = \frac{B_\text{DSB}}{2}.
\end{equation}

Para uma mensagem $m(t) = \cos(2\pi f_m t)$ e portadora $\cos(2\pi f_c t)$,
o produto vale
\begin{equation}
  s_\text{DSB}(t) = \tfrac{1}{2}\cos\!\left[2\pi (f_c - f_m) t\right]
                  + \tfrac{1}{2}\cos\!\left[2\pi (f_c + f_m) t\right],
\end{equation}
de modo que a componente em $f_c + f_m$ é a banda lateral superior, a
componente em $f_c - f_m$ é a banda lateral inferior e não existe componente
em $f_c$.

O sinal SSB também pode ser escrito pela representação de Hilbert:
\begin{align}
  s_\text{USB}(t) &= m(t)\cos(2\pi f_c t) - \hat{m}(t)\sin(2\pi f_c t), \\
  s_\text{LSB}(t) &= m(t)\cos(2\pi f_c t) + \hat{m}(t)\sin(2\pi f_c t),
\end{align}
onde $\hat{m}(t)$ é a transformada de Hilbert de $m(t)$, ou seja, uma
defasagem de $-90^\circ$ em todas as frequências positivas.

Na abordagem por filtragem, gera-se o DSB-SC e aplica-se um filtro passa-banda
que seleciona apenas a banda lateral superior ou a inferior. Nas duas
abordagens a demodulação é coerente: multiplica-se pelo oscilador local e
aplica-se um filtro passa-baixas.

Demodulando o sinal SSB-USB com um oscilador local $2\cos[2\pi(f_c+\Delta f)t]$
e filtrando as componentes em torno de $2f_c$, obtém-se
\begin{equation}
  y(t) = \mathrm{Re}\!\left\{ m_+(t)\, e^{-j 2\pi \Delta f t} \right\},
  \qquad m_+(t) = m(t) + j\hat{m}(t),
\end{equation}
isto é, o espectro da mensagem sofre uma translação rígida de $-\Delta f$.
Essa é a diferença essencial em relação ao DSB-SC, no qual o mesmo erro
produz duas cópias deslocadas de $\pm\Delta f$.

% ======================================================
\section{Materiais e blocos GNU Radio}

\begin{table}[H]
  \centering
  \begin{tabular}{l l l}
    \toprule
    \textbf{Bloco GRC} & \textbf{Localização} & \textbf{Função} \\
    \midrule
    \texttt{Variable}            & Miscellaneous        & Parâmetros globais \\
    \texttt{QT GUI Range} (2×)   & Instrumentation      & Sliders de $f_m$ e $\Delta f$ \\
    \texttt{Signal Source} (4×)  & Waveform Generators  & Mensagem, portadoras e oscilador local \\
    \texttt{Multiply}            & Math Operators       & Modulação e demodulação \\
    \texttt{Band Pass Filter}    & Filters              & Selecionar USB ou LSB \\
    \texttt{Low Pass Filter}     & Filters              & Recuperar a mensagem \\
    \texttt{Hilbert}             & Filters              & Sinal analítico da mensagem \\
    \texttt{Complex To Float}    & Type Converters      & Separar $m(t-\tau)$ e $\hat{m}(t-\tau)$ \\
    \texttt{Add} / \texttt{Subtract} & Math Operators   & Combinar os ramos I e Q \\
    \texttt{Delay}               & Misc Operators       & Alinhar mensagem e sinal demodulado \\
    \texttt{Throttle}            & Misc Operators       & Controle de CPU \\
    \texttt{QT GUI Time Sink}    & Instrumentation      & Formas de onda \\
    \texttt{QT GUI Freq Sink}    & Instrumentation      & Espectros \\
    \bottomrule
  \end{tabular}
  \caption{Blocos utilizados na Prática 04.}
\end{table}

% ======================================================
\section{Roteiro de montagem e execução}

\subsection*{Parâmetros iniciais}

\begin{center}
\begin{tabular}{l l}
  \toprule
  \texttt{samp\_rate} = 64\,000\,Hz    & \texttt{fc} = 16\,000\,Hz \\
  \texttt{fm} = 3\,000\,Hz (slider de 1\,500 a 6\,000\,Hz) & \texttt{delta\_f} = 0\,Hz (slider de $-500$ a $500$\,Hz) \\
  \texttt{bw\_trans} = 800\,Hz (banda de transição dos BPF) & \texttt{bw\_trans\_lp} = 600\,Hz (banda de transição do LPF) \\
  \texttt{n\_hilb} = 101 coeficientes  & \texttt{d\_align} = 224 amostras \\
  \bottomrule
\end{tabular}
\end{center}

Com esses valores, as componentes esperadas do DSB-SC ficam em
$f_c - f_m = 13\,000$\,Hz e $f_c + f_m = 19\,000$\,Hz.

\subsection*{Etapa 1 -- Gerar DSB-SC e identificar as bandas laterais (0--10 min)}

\begin{instrucao}[Geração e análise do DSB-SC]
\begin{enumerate}[label=\arabic*.]
  \item Crie as variáveis \texttt{samp\_rate = 64000} e \texttt{fc = 16000},
    e um \texttt{QT GUI Range} chamado \texttt{fm} com valor inicial 3\,000\,Hz,
    mínimo 1\,500\,Hz, máximo 6\,000\,Hz e passo 500\,Hz.
  \item Adicione o \texttt{Signal Source} da mensagem: Cosine, freq = \texttt{fm},
    amplitude 1,0.
  \item Adicione o \texttt{Signal Source} da portadora em fase: Cosine,
    freq = \texttt{fc}, amplitude 1,0.
  \item Adicione um \texttt{Multiply} e conecte as duas fontes.
  \item Conecte a saída ao \texttt{Throttle} e em seguida ao
    \texttt{QT GUI Freq Sink}, no canal 1.
  \item Execute e localize no espectro:
    \begin{itemize}
      \item banda lateral superior: componente em $f_c + f_m = 19\,000$\,Hz;
      \item banda lateral inferior: componente em $f_c - f_m = 13\,000$\,Hz;
      \item ausência de componente em $f_c = 16\,000$\,Hz.
    \end{itemize}
  \item Meça a largura de banda do DSB-SC como a distância entre as duas
    componentes e registre o nível de pico em dB.
  \item Mova o slider \texttt{fm} e confirme que as duas componentes se afastam
    simetricamente de $f_c$.
\end{enumerate}
\end{instrucao}

\subsection*{Etapa 2 -- Filtragem para obter a banda lateral superior (10--18 min)}

\begin{instrucao}[Selecionar a banda lateral superior]
\begin{enumerate}[label=\arabic*.]
  \item Adicione um \texttt{Band Pass Filter} do tipo \texttt{Float$\to$Float}
    após o \texttt{Throttle}:
    \begin{itemize}
      \item Decimation = 1, Gain = 1
      \item Low Cutoff Freq = \texttt{fc+fm/2}, ou seja 17\,500\,Hz
      \item High Cutoff Freq = \texttt{fc+2*fm}, ou seja 22\,000\,Hz
      \item Transition Width = \texttt{bw\_trans} = 800\,Hz
      \item Window = Hamming
    \end{itemize}
  \item Conecte a saída do filtro ao canal 2 da mesma \texttt{QT GUI Freq Sink}.
  \item Ajuste a escala vertical da Freq Sink para o intervalo de $-120$ a
    $10$\,dB, de modo que o resíduo da banda suprimida fique visível.
  \item Execute e verifique que resta apenas a componente em 19\,000\,Hz.
  \item Registre em dB o nível do resíduo em 13\,000\,Hz. Esse valor é a
    supressão da banda lateral indesejada obtida pelo método de filtragem.
\end{enumerate}
\end{instrucao}

\subsection*{Etapa 3 -- Filtragem para obter a banda lateral inferior (18--24 min)}

\begin{instrucao}[Selecionar a banda lateral inferior]
\begin{enumerate}[label=\arabic*.]
  \item Adicione um segundo \texttt{Band Pass Filter} em paralelo com o
    primeiro, também ligado à saída do \texttt{Throttle}:
    \begin{itemize}
      \item Low Cutoff Freq = \texttt{fc-2*fm}, ou seja 10\,000\,Hz
      \item High Cutoff Freq = \texttt{fc-fm/2}, ou seja 14\,500\,Hz
      \item Transition Width = \texttt{bw\_trans} = 800\,Hz
    \end{itemize}
  \item Conecte a saída ao canal 3 da \texttt{QT GUI Freq Sink}.
  \item Execute e verifique que resta apenas a componente em 13\,000\,Hz.
  \item Compare as três curvas e preencha as duas primeiras linhas da
    Tabela~\ref{tab:comp}.
\end{enumerate}
\end{instrucao}

\begin{atencao}[Por que as frequências de corte são escritas em função de \texttt{fm}]
As frequências de corte foram definidas como expressões e não como números
fixos. Assim, ao mover o slider \texttt{fm}, o GNU Radio recalcula os
coeficientes dos filtros e o experimento continua válido em toda a faixa do
slider. A banda de transição de 800\,Hz é menor do que a menor separação entre
as componentes, que é $2f_m = 3\,000$\,Hz para $f_m = 1\,500$\,Hz.
\end{atencao}

\subsection*{Etapa 4 -- Demodulação coerente do SSB (24--34 min)}

\begin{instrucao}[Demodulação coerente da banda lateral superior]
\begin{enumerate}[label=\arabic*.]
  \item Crie um \texttt{QT GUI Range} chamado \texttt{delta\_f}, com valor
    inicial 0\,Hz, mínimo $-500$\,Hz, máximo $500$\,Hz e passo 25\,Hz.
  \item Adicione o \texttt{Signal Source} do oscilador local: Cosine,
    freq = \texttt{fc+delta\_f}, amplitude 2,0.
  \item Adicione um \texttt{Multiply} e conecte a saída do filtro da banda
    lateral superior e o oscilador local.
  \item Adicione um \texttt{Low Pass Filter}:
    \begin{itemize}
      \item Cutoff Freq = \texttt{fm*1.5}
      \item Transition Width = \texttt{bw\_trans\_lp} = 600\,Hz
      \item Window = Hamming
    \end{itemize}
  \item Adicione um bloco \texttt{Delay} do tipo \texttt{float} com
    \texttt{Delay = d\_align = 224} e ligue nele a saída da fonte da mensagem.
  \item Conecte ao \texttt{QT GUI Time Sink}: canal 1 = saída do \texttt{Delay},
    canal 2 = saída do \texttt{Low Pass Filter}. Configure a Time Sink com
    \texttt{Autoscale = No}, \texttt{Y min} $=-1{,}2$, \texttt{Y max} $=1{,}2$ e
    \texttt{Number of Points = 256}.
  \item Execute com \texttt{delta\_f = 0} e verifique que as duas curvas ficam
    em fase. Meça a amplitude do sinal demodulado e compare com a amplitude
    unitária da mensagem.
\end{enumerate}
\end{instrucao}

\begin{atencao}[O bloco \texttt{Delay} e o valor 224]
O caminho de demodulação atravessa dois filtros FIR de fase linear. Com
janela de Hamming, o GNU Radio usa 193 coeficientes no passa-banda e 257
coeficientes no passa-baixas, o que corresponde a atrasos de grupo de 96 e 128
amostras. O total de 224 amostras depende apenas das bandas de transição, e
por isso permanece válido em toda a faixa do slider \texttt{fm}. Sem esse
bloco, as duas curvas da Time Sink aparecem deslocadas horizontalmente e a
comparação de amplitude fica prejudicada. Deixar \texttt{Autoscale} ligado
também mascara o resultado, porque a escala vertical é reajustada a cada
curva e o fator de escala deixa de ser visível.
\end{atencao}

\subsection*{Etapa 4b -- Erro de frequência do oscilador local (34--40 min)}

\begin{instrucao}[Medida do efeito de $\Delta f$]
\begin{enumerate}[label=\arabic*.]
  \item Conecte a saída do \texttt{Low Pass Filter} a uma segunda
    \texttt{QT GUI Freq Sink}.
  \item Com \texttt{delta\_f = 0}, confirme que o pico do sinal demodulado
    está em $f_m = 3\,000$\,Hz.
  \item Ajuste \texttt{delta\_f} para 100\,Hz e depois para 300\,Hz. Registre
    a frequência do pico em cada caso.
  \item Repita com \texttt{delta\_f} negativo e observe o sentido do
    deslocamento.
  \item Registre os valores medidos na Tabela~\ref{tab:df}.
\end{enumerate}
\end{instrucao}

\subsection*{Etapa 5 -- Geração SSB pela transformada de Hilbert (40--50 min)}

\begin{instrucao}[Método do defasamento em quadratura]
\begin{enumerate}[label=\arabic*.]
  \item Adicione um terceiro \texttt{Signal Source} para a portadora em
    quadratura: Sine, freq = \texttt{fc}, amplitude 1,0.
  \item Ligue a saída da fonte da mensagem a um bloco \texttt{Hilbert} com
    \texttt{Num Taps = n\_hilb = 101} e janela de Hamming.
  \item Ligue a saída do \texttt{Hilbert} a um bloco
    \texttt{Complex To Float}. A saída \texttt{re} é a mensagem atrasada,
    $m(t-\tau)$, e a saída \texttt{im} é a transformada de Hilbert,
    $\hat{m}(t-\tau)$, com $\tau = (\texttt{n\_hilb}-1)/2 = 50$ amostras.
  \item Monte os dois ramos:
    \begin{itemize}
      \item ramo I: saída \texttt{re} multiplicada pela portadora em fase;
      \item ramo Q: saída \texttt{im} multiplicada pela portadora em quadratura.
    \end{itemize}
  \item Combine os ramos: um bloco \texttt{Subtract} calculando
    $\text{I} - \text{Q}$ produz a banda lateral superior e um bloco
    \texttt{Add} calculando $\text{I} + \text{Q}$ produz a banda lateral
    inferior.
  \item Conecte as duas saídas a uma terceira \texttt{QT GUI Freq Sink}, com
    escala vertical de $-120$ a $10$\,dB.
  \item Execute e registre o nível do resíduo da banda suprimida, do mesmo modo
    que na Etapa 2. Compare com o valor obtido pelo método de filtragem.
\end{enumerate}
\end{instrucao}

\begin{atencao}[{O ramo em fase sai do bloco \texttt{Hilbert}, não da fonte}]
O filtro de Hilbert é um FIR de fase linear e atrasa o sinal em
$(\texttt{n\_hilb}-1)/2$ amostras. O bloco \texttt{Hilbert} do GNU Radio já
devolve na parte real a mensagem atrasada desse mesmo valor, de forma que os
dois ramos saem alinhados. Se o ramo I for ligado diretamente à fonte da
mensagem, os dois ramos ficam desalinhados no tempo, o cancelamento falha e as
duas bandas laterais continuam presentes na saída. Uma alternativa equivalente
é usar a mensagem original no ramo I e inserir um bloco \texttt{Delay} de 50
amostras antes do multiplicador.
\end{atencao}

% ======================================================
\section{Tabelas de registro}

\begin{table}[H]
  \centering
  \begin{tabular}{l c c c}
    \toprule
    \textbf{Parâmetro} & \textbf{DSB-SC} & \textbf{SSB-USB} & \textbf{SSB-LSB} \\
    \midrule
    Largura de banda (Hz)                    & & & \\
    Número de componentes espectrais         & & & \\
    Frequência das componentes (Hz)          & & & \\
    Nível da componente principal (dB)       & & & \\
    Resíduo da banda suprimida por filtragem (dB) & --- & & \\
    Resíduo da banda suprimida por Hilbert (dB)   & --- & & \\
    \bottomrule
  \end{tabular}
  \caption{Comparação entre DSB-SC, SSB-USB e SSB-LSB. Conte apenas as
    componentes em frequências positivas, que são as exibidas pela
    \texttt{QT GUI Freq Sink}.}
  \label{tab:comp}
\end{table}

\begin{table}[H]
  \centering
  \begin{tabular}{c c c}
    \toprule
    $\Delta f$ (Hz) & Frequência medida do pico (Hz) & Valor previsto $f_m - \Delta f$ (Hz) \\
    \midrule
    $0$    & & \\
    $100$  & & \\
    $300$  & & \\
    $-200$ & & \\
    \bottomrule
  \end{tabular}
  \caption{Efeito do erro de frequência do oscilador local sobre o sinal
    demodulado.}
  \label{tab:df}
\end{table}

% ======================================================
\section{Perguntas para o relatório}

\begin{enumerate}[label=\textbf{Q\arabic*.}, leftmargin=*]

  \item Como você identificou a banda lateral superior e a inferior no
    espectro do DSB-SC? Que característica do espectro permite diferenciá-las
    e por que não há componente em $f_c$?

  \item Qual foi a largura de banda medida do sinal SSB em comparação com a do
    DSB-SC? Confirme a relação $B_\text{SSB} = B_\text{DSB}/2$ com os dados da
    Tabela~\ref{tab:comp} para pelo menos dois valores de $f_m$.

  \item A amplitude do sinal demodulado medida na Etapa 4 é diferente da
    amplitude da mensagem original. Deduza analiticamente o fator de escala e
    compare com o valor medido. Em seguida, cite duas vantagens e uma
    desvantagem práticas do SSB em relação ao DSB-SC.

  \item Use a Tabela~\ref{tab:df} para responder: o que acontece com o
    espectro da mensagem recuperada quando o oscilador local tem desvio
    $\Delta f$? Mostre analiticamente que no SSB o espectro sofre uma
    translação rígida de $-\Delta f$, enquanto no DSB-SC o mesmo erro produz
    duas cópias deslocadas de $\pm\Delta f$. Explique por que a translação
    rígida é mais danosa para um sinal de voz.

  \item Compare os dois métodos de geração usando os resíduos medidos nas
    Etapas 2 e 5. Qual apresentou melhor supressão da banda lateral
    indesejada? Qual é mais simples de implementar em hardware analógico?
    Discuta o que acontece com cada método quando a mensagem contém
    componentes de frequência muito próximas de zero.

\end{enumerate}

% ======================================================
\section{Entregáveis}

O entregável é o relatório no modelo definido no SIGAA contendo:

\begin{itemize}
  \item Captura do espectro mostrando DSB-SC, banda lateral superior e banda
    lateral inferior obtidas por filtragem.
  \item Captura da comparação temporal entre a mensagem original e a
    demodulada, com \texttt{Autoscale} desligado.
  \item Captura do espectro demodulado para pelo menos dois valores de
    \texttt{delta\_f}.
  \item Captura do espectro das bandas laterais geradas pelo método de Hilbert.
  \item Tabelas~\ref{tab:comp} e \ref{tab:df} preenchidas.
  \item Respostas às 5 questões.
\end{itemize}

\begin{atencao}[Sobre os parâmetros dos filtros]
Ajuste as frequências de corte com cuidado: a banda de transição deve ser
menor do que a separação entre as componentes espectrais, que vale $2f_m$.
Se o filtro passa-banda não atenuar suficientemente a banda indesejada,
reduza a banda de transição, o que aumenta automaticamente o número de
coeficientes, ou troque a janela de Hamming pela de Kaiser. No método de
Hilbert, aumentar \texttt{n\_hilb} melhora a supressão sobretudo nas
frequências mais baixas da mensagem.
\end{atencao}

\end{document}
